% OakGPT Chat C2 report. LaTeX source of report.md; the two must stay in sync.
% Build (pdfLaTeX, XeLaTeX or Tectonic) from this directory, e.g. `tectonic report.tex`.
% The source is pure ASCII so every engine accepts it.
\documentclass[11pt,a4paper]{article}

\usepackage{iftex}
\ifPDFTeX
  \usepackage[T1]{fontenc}
  \usepackage{lmodern}
\fi
\usepackage[margin=25mm]{geometry}
\usepackage{microtype}
\usepackage{booktabs}
\usepackage{tabularx}
\usepackage{array}
\usepackage{xcolor}
\usepackage{enumitem}
\usepackage{fancyvrb}
\usepackage{caption}
\usepackage{float}
\usepackage[hidelinks]{hyperref}

\definecolor{oak}{HTML}{3D6B35}
\hypersetup{colorlinks=true,linkcolor=oak,urlcolor=oak,citecolor=oak,
  pdftitle={OakGPT Chat C2: LoRA Fine-Tuning on a 6 GB GPU},
  pdfauthor={OakGPT - Local LLM}}
\setlist{itemsep=2pt,topsep=4pt}
\setlength{\parindent}{0pt}
\setlength{\parskip}{6pt}
\newcolumntype{Y}{>{\raggedright\arraybackslash}X}
\newcolumntype{R}{>{\raggedleft\arraybackslash}X}
\renewcommand{\arraystretch}{1.15}
\urlstyle{same}
\setlength{\emergencystretch}{3em}

\newcommand{\code}[1]{\texttt{#1}}
\title{\textbf{OakGPT Chat C2:\\LoRA Fine-Tuning on a 6\,GB GPU}\\[4pt]
  \large Design decisions, rationale, sources and measured results}
\author{OakGPT -- Local LLM}
\date{2026-11-26}

\begin{document}
\maketitle



\begin{table}[H]\centering\small
\begin{tabularx}{\textwidth}{@{}lY@{}}
\toprule
Track & Chat C2, fine-tuning as a measured hypothesis \\
Protocol & protocol.md, frozen in commit \code{215881b} before training \\
Training data & \code{data/finetune/c2\_distill\_v1}: 1,076 prompts, 970 accepted teacher answers \\
Evaluation & \code{chat\_eval\_v1} text items (40), frozen in \code{851c057}; never used for training \\
Raw results & raw/: every reply, training log, memory probes, hashes, environment \\
Hardware & NVIDIA GeForce RTX 2060, 6 GB VRAM (about 5.0 GB free), Windows 10 19045 \\
Training stack & torch 2.9.1+cu126, transformers 4.57.1, bitsandbytes 0.49.2, accelerate 1.11.0 \\
Inference & llama.cpp \code{b11930}, Qwen3 1.7B Q4\_K\_M + LoRA GGUF (\code{--lora}) \\
Adapter & \code{c2-qwen3-1.7b-lora-v1-f16.gguf}, SHA-256 \code{2ead4472...707b0d} (not committed) \\
\bottomrule\end{tabularx}\end{table}

\section{Summary}

\begin{itemize}
\item \textbf{Feasibility: yes.} QLoRA training of Qwen3 1.7B fits comfortably on the RTX 2060: 1.7--3.2 GB peak during training steps with gradient checkpointing. Without checkpointing it does not fit. The whole C2 adapter trained in 6.6 minutes.
\item \textbf{Fast mode: large gain.} Strict accuracy rose from 49.2\% to 74.2\% (+25.0 points, 95\% CI 12.5 to 38.3; 12 items better, 1 worse).
\item \textbf{Reasoning mode: large loss.} The adapter is also active when thinking is on, and it taught the model to stop thinking. Reasoning-mode strict accuracy fell from 91.7\% to 56.7\% ($-$35.0 points, CI $-$50.0 to $-$20.0).
\item \textbf{Decision (pre-registered rule): do not adopt.} Criterion 2 fails, so the app keeps the unchanged base model. The result points to a concrete follow-up: apply the adapter only to fast-mode requests (llama-server supports a per-request LoRA scale), evaluated on a new frozen set.
\end{itemize}

\section{What was built}

\begin{itemize}
\item \code{src/oakgpt/finetuning/lora.py}: a small explicit LoRA implementation (Hu et al. 2021 [1]). It wraps linear layers and saves the adapter with PEFT tensor names, which llama.cpp's \code{convert\_lora\_to\_gguf.py} reads. CPU tests cover identity at initialization, the update formula, trainable-parameter isolation and a save/load round trip.
\item \code{src/oakgpt/finetuning/train.py}: QLoRA training, memory probes and WikiText perplexity in the separate \code{.venv-finetune} environment (\code{scripts/setup\_finetune.ps1}).
\item \code{src/oakgpt/finetuning/distill.py} and \code{oakgpt c2-distill}: teacher answers through the app's own code path; \code{oakgpt c2-open-probe}: fast-mode reply lengths.
\item \code{scripts/build\_c2\_prompts\_v1.py}: the deterministic prompt set with a contamination guard; a test rebuilds it byte for byte.
\item \code{managed\_server.lora} in the chat configuration, which passes \code{--lora} to llama-server and records the adapter hash in the status identity.
\end{itemize}

\section{Decision records}

\subsection{D33. Write LoRA ourselves; use transformers and bitsandbytes for the base}

\textbf{Decision.} Implement LoRA itself in about 150 lines. Load the base model with transformers and quantize it with bitsandbytes. Keep all of it in a separate optional environment.

\textbf{Rationale.}

\begin{itemize}
\item LoRA is a small, well-specified operation. Writing it follows the project rule to prefer small explicit implementations, and makes it testable on CPU without downloads.
\item Loading Qwen3 needs its model code (transformers 4.51 or later). Writing our own NF4 kernels would be neither small nor trustworthy.
\item The core installation stays \code{torch} only. The existing \code{.venv-cuda} (Phase 2) is untouched.
\end{itemize}

\textbf{Alternatives.}

\begin{itemize}
\item PEFT is the standard library for this. It is rejected only because the required part is tiny. Our files use its format, so PEFT-based tools still work.
\item Unsloth and Axolotl are larger frameworks than this experiment needs.
\end{itemize}

\subsection{D34. QLoRA (NF4 base) rather than fp16 LoRA}

\textbf{Decision.} Train on a 4-bit NF4 base with double quantization and fp16 compute, as in QLoRA (Dettmers et al. 2023 [2]).

\textbf{Evidence.}

\begin{itemize}
\item fp16 LoRA also ran, at 4.2 GB peak against 5.0 GB free. That leaves about 0.3 GB of margin with an otherwise idle desktop.
\item NF4 peaked at 1.7--2.2 GB for 512-token sequences.
\item The speed was similar: 578--910 tokens/s for either.
\end{itemize}

\textbf{Rationale.} The memory margin decides. A browser or a game on the same GPU (see Limitations) would break fp16 LoRA. Turing GPUs have no bf16, so compute is fp16 with fp32 adapter weights and loss scaling.

\subsection{D35. Gradient checkpointing is required; measure memory, do not trust ``it ran''}

\textbf{Decision.} Always train with gradient checkpointing. Count a configuration as fitting only if its measured peak stays within physical VRAM.

\textbf{Evidence.} Without checkpointing, both probed NF4 configurations failed. Before failing, PyTorch reported 12.6 GB allocated on a 6 GB card. The Windows driver can let allocations spill into shared system memory, so a configuration can appear to run while actually using system RAM, and run very slowly. The probes therefore record peak allocation and device capacity, not just success.

A first probe series was invalid: the model was still in evaluation mode, where checkpointing is inactive, and its peak included model loading. Those files are kept in \code{raw/feasibility-v1-loadpeak/}; the table below is from the corrected series.

\begin{table}[H]\centering\small
\begin{tabularx}{\textwidth}{@{}Ylll@{}}
\toprule
Configuration (seq $\times$ micro-batch) & Peak allocated & Median step & Tokens/s \\
\midrule
NF4, 512 $\times$ 1, checkpointing & 1.66 GB & 0.89 s & 578 \\
NF4, 512 $\times$ 4, checkpointing & 2.20 GB & 3.09 s & 662 \\
NF4, 1024 $\times$ 4, checkpointing & 3.21 GB & 5.09 s & 805 \\
NF4, 512 $\times$ 4, no checkpointing & out of memory (12.65 GB requested) & -- & -- \\
NF4, 256 $\times$ 8, no checkpointing & out of memory (12.55 GB requested) & -- & -- \\
fp16, 512 $\times$ 1, checkpointing & 3.61 GB & 0.63 s & 810 \\
fp16, 512 $\times$ 4, checkpointing & 4.15 GB & 2.25 s & 910 \\
\bottomrule\end{tabularx}\end{table}

Notes on the table:

\begin{itemize}
\item Peaks cover training steps only. Loading peaked separately, at 2.42 GB for NF4 and 4.36 GB for fp16.
\item The inputs are synthetic tokens with 64 labelled positions per row and LoRA r=16 on all projections (17.4 M trainable parameters).
\item Step times vary between runs: the same NF4 512$\times$4 setting took 2.25 s in the first series.
\end{itemize}

\subsection{D36. Self-distillation: the model's own reasoning answers train its fast mode}

\textbf{Decision.} Generate new prompts and let the same model answer them in reasoning mode. Keep a keyed answer only if it is strictly correct, and train fast mode on the final answers, never on the thinking text.

\textbf{Rationale.}

\begin{itemize}
\item Evaluation v1 showed that reasoning mode already knows most answers (92\%), but costs about 2 s and up to 63 s per reply.
\item Moving that behaviour into the cheap mode is the idea of ``Distilling System 2 into System 1'' (Yu et al. 2024 [3]). Filtering by correctness is rejection sampling as in STaR (Zelikman et al. 2022 [4]).
\item The teacher is local and the same model, so nothing leaves the computer and no outside model's text or licence enters the data.
\end{itemize}

\textbf{Result.}

\begin{itemize}
\item Collection took 26 minutes. 970 of 1,076 answers were accepted: 882 for training and 88 for a development split.
\item Rejections were genuine errors. Most were in Swedish (47 of 125, for example ``Copenhagen'' for ``K\"openhamn''), instructions (27, for example \code{lxeip} for ``pixel'' backwards) and code (24, for example \code{2,6} for \code{print(24 // 9, 24 \% 9)}).
\item 77 open-ended prompts (explanations, tips, short messages) were included so the adapter would not learn only one-word replies.
\end{itemize}

\subsection{D37. A contamination guard against the frozen evaluation set}

\textbf{Decision.} The prompt builder rejects a generated prompt if any of these hold:

\begin{itemize}
\item its normalized text equals an evaluation prompt
\item its content without the answer-format sentence equals an evaluation item's content
\item it contains a term reserved for an evaluation item (for example ``Australia'', ``\"apple'', ``M\"alardalen'', ``Pepper'')
\item it shares any 8-word sequence with an evaluation prompt
\end{itemize}

\textbf{Result.}

\begin{itemize}
\item 139 prompts were rejected (2 exact, 137 by 8-grams) and 73 duplicates were dropped.
\item Whole templates disappeared this way, for example the weekday and discount word problems and the ``meeting moved'' reading texts, because their phrasing followed the evaluation items. \code{contamination.json} lists the rejected templates.
\item A test checks that the committed set has no evaluation prompt or content.
\end{itemize}

\textbf{Caveat.} The generator was written knowing the evaluation's skill categories. Gains therefore show transfer to new items of the same kinds, not a general improvement. Model answers on generated data can still coincide with evaluation answers (for example the number 13).

\subsection{D38. Serve the adapter in llama.cpp on the unchanged Q4\_K\_M model}

\textbf{Decision.} Convert the adapter with llama.cpp's own \code{convert\_lora\_to\_gguf.py} (f16 output, 35 MB) and load it in llama-server with \code{--lora}. Do not merge it and requantize.

\textbf{Rationale.} The app's serving path, quantization, prompt cache and model swapping stay identical. Only the adapter is added, and it can be removed again.

\textbf{Check.} The adapter was trained against an NF4 base but is applied to Q4\_K\_M, a mismatch that is measured rather than assumed harmless. An all-zero adapter reproduced the base model's replies exactly on 12 development prompts with the same seed. The trained adapter changed 8 of them (\code{raw/lora-check.json}).

\subsection{D39. Pre-registration, same-session baseline, one training run}

\textbf{Decision.} Freeze the protocol, data, hyperparameters and decision rule in Git before training, and run one training with the recipe's defaults:

\begin{itemize}
\item r=16, alpha=32, dropout 0.05 on all 7 projections
\item lr 2e-4 with cosine decay, 2 epochs, effective batch 16
\end{itemize}

The final adapter is used, with no checkpoint selection. The base model is re-measured in the same session.

\textbf{Rationale.}

\begin{itemize}
\item Tuning on the evaluation set would turn it into a training signal.
\item A same-session baseline separates the adapter's effect from drift.
\end{itemize}

\textbf{Result.} The re-measured base reproduced the v1 run exactly: 0 of 40 items differed in either mode. So the comparison below is not affected by run-to-run noise.

\subsection{D40. Outcome: keep the base model; scope the adapter to fast mode next}

\textbf{Decision.} Do not adopt the adapter. \code{configs/chat/qwen3\_q4km\_c2.json} remains as an experiment configuration only. The app's default stays reasoning mode on the base model.

\textbf{Rationale.} The pre-registered criterion 2 failed by a wide margin. See Interpretation for the follow-up.

\section{Results}

\subsection{Training}

\begin{itemize}
\item 882 examples, 112 optimizer steps in 397 s. The median step took 3.2 s for a median of 141 answer tokens.
\item Development loss on answer tokens fell from 1.915 to 0.784 after epoch 1 and 0.780 after epoch 2.
\item 17.4 M parameters were trainable, 1.0\% of the model.
\end{itemize}

\subsection{Evaluation on \code{chat\_eval\_v1} (40 text items $\times$ 3 seeds)}

\begin{table}[H]\centering\small
\begin{tabularx}{\textwidth}{@{}Ylll@{}}
\toprule
Condition & Base & Adapter & Paired difference (95\% CI) \\
\midrule
Fast mode, strict & 49.2\% & \textbf{74.2\%} & \textbf{+25.0} (+12.5 to +38.3) \\
Fast mode, lenient & 60.8\% & 74.2\% &  \\
Reasoning mode, strict & 91.7\% & 56.7\% & \textbf{$-$35.0} ($-$50.0 to $-$20.0) \\
Fast mode, median wall time & 0.065 s & 0.078 s &  \\
Reasoning mode, median wall time & 2.13 s & 0.16 s &  \\
Reasoning mode, median generated tokens & 155 & 8 &  \\
\bottomrule\end{tabularx}\end{table}

Per-category strict accuracy:

\begin{table}[H]\centering\small
\begin{tabularx}{\textwidth}{@{}Yllll@{}}
\toprule
Category (items) & Fast, base & Fast, adapter & Reasoning, base & Reasoning, adapter \\
\midrule
Arithmetic (6) & 83\% & 83\% & 100\% & 67\% \\
Word problems (6) & 50\% & 50\% & 89\% & 89\% \\
Facts (6) & 83\% & 100\% & 100\% & 11\% \\
Swedish (6) & 6\% & 39\% & 83\% & 22\% \\
Python code (7) & 24\% & 86\% & 86\% & 76\% \\
Instructions (5) & 73\% & 67\% & 87\% & 47\% \\
Reading (4) & 25\% & 100\% & 100\% & 92\% \\
\bottomrule\end{tabularx}\end{table}

\subsection{Other checks}

\begin{itemize}
\item \textbf{Open-ended fast-mode replies} on 12 prompts outside the training data had a median of 108.5 tokens for the base and 113.5 for the adapter, a ratio of 1.05. Replies stayed full explanations; the adapter did not make the model terse. One adapter reply hit the token limit.
\item \textbf{WikiText-2 test}, 64 windows of 1,024 tokens on the NF4 base: NLL per token was 2.957 for the base and 2.792 for the adapter (perplexity 19.2 versus 16.3). The adapter did not degrade general text modelling on this measure. The improvement is unexpected and not explained here; it is not used as evidence of better quality.
\end{itemize}

\subsection{Decision rule}

\begin{table}[H]\centering\small
\begin{tabularx}{\textwidth}{@{}Yll@{}}
\toprule
Criterion & Result & Pass \\
\midrule
1. Fast-mode gain with a CI that excludes zero & +25.0 (12.5 to 38.3) & yes \\
2. Reasoning drop of at most 5 points & $-$35.0 & \textbf{no} \\
3. WikiText NLL rise of at most 5\% & $-$5.6\% & yes \\
4. Open-ended reply length at least half of base & 1.05$\times$ & yes \\
\bottomrule\end{tabularx}\end{table}

\section{Interpretation}

\textbf{Why fast mode improved.} The biggest gains are where evaluation v1 found format or knowledge failures:

\begin{itemize}
\item Reading went from 25\% to 100\% and code from 24\% to 86\%. The base model knew these answers but replied with sentences or extra text; the adapter answers in the requested form, which is what self-distillation targets. Strict and lenient accuracy are now equal (74.2\%).
\item Swedish rose from 6\% to 39\%, for words that were not in the training data. Swedish vocabulary remains the weakest area.
\item Word problems did not improve (50\%). Multi-step arithmetic still needs written-out steps, as Yu et al. [3] also observed for chain-of-thought tasks.
\end{itemize}

\textbf{Why reasoning mode broke.} In reasoning mode the prompt ends right before the model would write \code{<think>}. Training only ever showed the model answering directly after an empty thinking block, and the adapter generalized that to ``answer immediately'':

\begin{itemize}
\item The median thinking text fell from 495 to 5 characters, and 38 of 120 replies had none.
\item Often the model wrote the answer where its thoughts would go and never closed the block, so the visible answer was empty. For example, for \code{1000 - 387} the thinking channel contained \code{713} and the reply was empty.
\item Facts collapsed to 11\% this way.
\end{itemize}

The loss is a side effect of applying one adapter to both modes, not of the knowledge the adapter learned.

\textbf{Next step.} llama-server accepts a per-request \code{lora} list with a scale per adapter (\code{[\{"id": 0, "scale": 1.0\}]}, server README, b11930). The app could send scale 1 in fast mode and 0 in reasoning mode, without restarting or swapping models. Fast mode would then get the measured gain and reasoning mode the unchanged base.

That design was not part of this protocol, and \code{chat\_eval\_v1} has now been used to select it. It therefore needs its own pre-registered evaluation on a new frozen set (\code{chat\_eval\_v2}). It also needs a check of how prompt caching behaves when the scale changes between turns.

\section{Limitations}

\begin{itemize}
\item \textbf{Small evaluation.} 40 text items. The intervals are wide, and per-category numbers rest on 4--7 items each.
\item \textbf{In-distribution training.} Training covered the same skills as the evaluation (D37). The gain is not a general capability measure.
\item \textbf{One training run.} One seed, default hyperparameters, no variance across training runs.
\item \textbf{Quantization mismatch.} Trained on NF4, served on Q4\_K\_M. The zero-adapter check covers the serving path, not the mismatch itself.
\item \textbf{Shared GPU.} A game used the GPU during the first, discarded probe attempt. Measurements were paused until it closed and restarted on a free GPU. A watcher found no game running during the probes, distillation, training or evaluation. Desktop use still varies, so absolute memory margins depend on what else is running.
\item \textbf{WikiText} measures raw text, not chat quality. The open-ended probe measures length, not correctness.
\end{itemize}

\section{Reproduction}

\begin{Verbatim}[fontsize=\footnotesize]
powershell -NoProfile -File scripts/setup_finetune.ps1
.\.venv\Scripts\python scripts\build_c2_prompts_v1.py data\finetune\c2_distill_v1   # into an empty directory
.\.venv\Scripts\oakgpt c2-distill --config configs/chat/qwen3_q4km.json `
  --prompts data/finetune/c2_distill_v1/prompts.jsonl --output <new directory>
.\.venv-finetune\Scripts\python -m oakgpt.finetuning.train train `
  --config configs/finetune/c2_qlora_v1.json --output runs/c2/train-v1
.\.venv-convert\Scripts\python vendor/llama.cpp-source/convert_lora_to_gguf.py --base models/qwen3-original `
  --outtype f16 --outfile models/c2-qwen3-1.7b-lora-v1-f16.gguf runs/c2/train-v1/adapter
powershell -NoProfile -File scripts/c2_evaluate.ps1
.\.venv\Scripts\python scripts\c2_compare.py runs/c2 <new summary.json>
\end{Verbatim}

\code{scripts/c2\_feasibility.ps1} repeats the memory probes and \code{scripts/c2\_lora\_check.py} the zero-adapter check. All commands refuse to overwrite existing outputs.

\section*{References}

\begin{enumerate}
\item E. J. Hu et al., ``LoRA: Low-Rank Adaptation of Large Language Models,'' arXiv:2106.09685, 2021.
\item T. Dettmers, A. Pagnoni, A. Holtzman, L. Zettlemoyer, ``QLoRA: Efficient Finetuning of Quantized LLMs,'' NeurIPS 2023, arXiv:2305.14314.
\item P. Yu, J. Xu, J. Weston, I. Kulikov, ``Distilling System 2 into System 1,'' arXiv:2407.06023, 2024.
\item E. Zelikman, Y. Wu, J. Mu, N. D. Goodman, ``STaR: Bootstrapping Reasoning With Reasoning,'' NeurIPS 2022, arXiv:2203.14465.
\item T. Chen, B. Xu, C. Zhang, C. Guestrin, ``Training Deep Nets with Sublinear Memory Cost,'' arXiv:1604.06174, 2016 (gradient checkpointing).
\item llama.cpp server documentation, \code{tools/server/README.md} at tag b11930 (\code{--lora}, per-request \code{lora} field); \code{convert\_lora\_to\_gguf.py}.
\item Qwen Team, ``Qwen3-1.7B'' model card (thinking and non-thinking modes, chat template).
\end{enumerate}
\end{document}
